\section{Metodología de Evaluación}

Para evaluar los resultados y el impacto de la implementación de Odoo 19 en Machu Picchu
Exclusive Tours E.I.R.L., se aplicó un enfoque de análisis integral orientado a medir la
profesionalización del ciclo comercial y la seguridad operativa. La evaluación se estructuró
bajo los siguientes enfoques:

\begin{itemize}
    \item \textbf{Comparación de procesos (Antes vs. Después):} se realizó un análisis comparativo entre el modelo manual ---basado en WhatsApp y Excel--- frente al nuevo entorno centralizado en Odoo. Se evaluó especialmente la eliminación de la dispersión de datos y la mejora en la trazabilidad de los itinerarios de viaje, que varían de 2 a 20 días.
    \item \textbf{Análisis de eficiencia operativa:} se verificó el desempeño de los flujos de cotización y el registro de adelantos (del 30\% al 50\%) directamente en el sistema. El enfoque se centró en medir la reducción de tiempos de respuesta a clientes extranjeros y la eliminación de errores en la compra de boletos no reembolsables.
    \item \textbf{Evaluación de indicadores de desempeño (KPIs):} se contrastaron los indicadores definidos en el diagnóstico inicial con los resultados obtenidos tras la puesta en marcha. Se midieron métricas clave como el crecimiento de reservas digitales, la tasa de conversión de cotizaciones y la seguridad en el almacenamiento de pasaportes.
    \item \textbf{Observación directa y nivel de adopción:} se realizó un seguimiento del uso del CRM y el módulo de Proyectos en el entorno real de trabajo. Este análisis permitió evidenciar la reducción en el uso de herramientas externas (como el borrado manual de fotos de pasaportes en teléfonos personales) y el grado de autonomía alcanzado por el equipo.
\end{itemize}
